\documentclass[10pt,a4paper]{amsart}
\usepackage{amsmath,amssymb,amsthm,graphicx}
% Retained minimal preamble for the Stacks Project chapter more-morphisms.tex (More on Morphisms).
% Verbatim excerpt of the chapter preamble: Xy-pic setup (native preamble lines 14-18) and the native
% theorem declarations with the mathematical macros (native preamble lines 154-217). Copied mechanically;
% nothing added, removed or redefined. The wrapper class and the amsmath/amssymb/amsthm/graphicx base
% packages are supplied by the assembler, so the class/externaldocument part of the native preamble is omitted.
% For commutative diagrams we use Xy-pic
\usepackage[all]{xy}

% We use 2cell for 2-commutative diagrams.
\xyoption{2cell}
% Theorem environments.
%
\theoremstyle{plain}
\newtheorem{theorem}[subsection]{Theorem}
\newtheorem{proposition}[subsection]{Proposition}
\newtheorem{lemma}[subsection]{Lemma}

\theoremstyle{definition}
\newtheorem{definition}[subsection]{Definition}
\newtheorem{example}[subsection]{Example}
\newtheorem{exercise}[subsection]{Exercise}
\newtheorem{situation}[subsection]{Situation}

\theoremstyle{remark}
\newtheorem{remark}[subsection]{Remark}
\newtheorem{remarks}[subsection]{Remarks}

\numberwithin{equation}{subsection}

% Macros
%
\def\lim{\mathop{\mathrm{lim}}\nolimits}
\def\colim{\mathop{\mathrm{colim}}\nolimits}
\def\Spec{\mathop{\mathrm{Spec}}}
\def\Hom{\mathop{\mathrm{Hom}}\nolimits}
\def\Ext{\mathop{\mathrm{Ext}}\nolimits}
\def\SheafHom{\mathop{\mathcal{H}\!\mathit{om}}\nolimits}
\def\SheafExt{\mathop{\mathcal{E}\!\mathit{xt}}\nolimits}
\def\Sch{\mathit{Sch}}
\def\Mor{\mathop{\mathrm{Mor}}\nolimits}
\def\Ob{\mathop{\mathrm{Ob}}\nolimits}
\def\Sh{\mathop{\mathit{Sh}}\nolimits}
\def\NL{\mathop{N\!L}\nolimits}
\def\CH{\mathop{\mathrm{CH}}\nolimits}
\def\proetale{{pro\text{-}\acute{e}tale}}
\def\etale{{\acute{e}tale}}
\def\QCoh{\mathit{QCoh}}
\def\Ker{\mathop{\mathrm{Ker}}}
\def\Im{\mathop{\mathrm{Im}}}
\def\Coker{\mathop{\mathrm{Coker}}}
\def\Coim{\mathop{\mathrm{Coim}}}

% Boxtimes
%
\DeclareMathSymbol{\boxtimes}{\mathbin}{AMSa}{"02}

%
% Macros for moduli stacks/spaces
%
\def\QCohstack{\mathcal{QC}\!\mathit{oh}}
\def\Cohstack{\mathcal{C}\!\mathit{oh}}
\def\Spacesstack{\mathcal{S}\!\mathit{paces}}
\def\Quotfunctor{\mathrm{Quot}}
\def\Hilbfunctor{\mathrm{Hilb}}
\def\Curvesstack{\mathcal{C}\!\mathit{urves}}
\def\Polarizedstack{\mathcal{P}\!\mathit{olarized}}
\def\Complexesstack{\mathcal{C}\!\mathit{omplexes}}
% \Pic is the operator that assigns to X its picard group, usage \Pic(X)
% \Picardstack_{X/B} denotes the Picard stack of X over B
% \Picardfunctor_{X/B} denotes the Picard functor of X over B
\def\Pic{\mathop{\mathrm{Pic}}\nolimits}
\def\Picardstack{\mathcal{P}\!\mathit{ic}}
\def\Picardfunctor{\mathrm{Pic}}
\def\Deformationcategory{\mathcal{D}\!\mathit{ef}}
\usepackage{cleveref}
\crefname{originallemma}{original lemma}{original lemmas}
\Crefname{originallemma}{Original lemma}{Original lemmas}
\makeatletter
\newlabel{lemma-count-geometric-fibres}{{0F35}{source}{}{sourcepointer-1}{}}
\newlabel{lemma-count-geometric-fibres@cref}{{[originallemma][0][]0F35}{[1][0][]source}{}{}{}}
\makeatother
\crefname{originallemma}{original lemma}{original lemmas}
\Crefname{originallemma}{Original lemma}{Original lemmas}
\makeatletter
\newlabel{lemma-etale-splits-off-quasi-finite-part-technical-variant}{{02LO}{source}{}{sourcepointer-2}{}}
\newlabel{lemma-etale-splits-off-quasi-finite-part-technical-variant@cref}{{[originallemma][0][]02LO}{[1][0][]source}{}{}{}}
\makeatother
\title{Lower semicontinuity of the geometric fibre count for universally open morphisms}
\author{The Stacks Project Authors}
\date{stacks-project:more-morphisms.tex:local-snapshot-20261011; GNU FDL 1.2 or later}
\begin{document}
\maketitle
\noindent\textit{Editorial note.} Counted span is the author's own proof environment, source lines 21989-22016, immediately following the statement at 21980-21987.\par\medskip
\noindent\textit{Editorial note.} The author writes 'details omitted' twice, at source lines 21994 and 22008; both sentences are retained verbatim. The first shrinks to a quasi-compact open piece, the second identifies n\_\{X/Y\}(y) with the sum of the multiplicities after the etale splitting; the displayed computation that uses those two facts is inside the span.\par\medskip
\noindent\textit{Editorial note.} The statement itself uses \textbackslash{}ref to lemma-count-geometric-fibres to name the invariant; it is kept as a navigation pointer.\par\medskip
\noindent\textit{Allowed standard prerequisites (editorial).}\begin{itemize}
\item Definition of the count of geometric fibres n\_\{X/Y\}. chapter lemma lemma-count-geometric-fibres, cited at proof line 21984. Fixes the invariant whose lower semicontinuity is asserted and records that it is computed from the connected components of the geometric fibre.
\item Etale splitting off of the quasi-finite part. chapter lemma lemma-etale-splits-off-quasi-finite-part-technical-variant, cited at proof lines 21997-21998. Produces the etale neighbourhood and the decomposition of U x\_Y X into W and the finite list of V\_\{i,j\} which exhibits the count as a sum of multiplicities.
\item Base change compatibility of n\_\{X/Y\}. used at proof lines 22014-22015. Transfers the inequality proved over the etale neighbourhood back to Y.
\end{itemize}\medskip
\ifdefined\hypertarget\hypertarget{sourcepointer-1}{}\fi
\noindent\textit{Source reference.} Original lemma 0F35: sources/stacks-more-morphisms/pinned-original-source.tex, source lines 21960--21960. Navigation pointer only: the counted proof cites this statement by \textbackslash{}ref as a standing input; the deduction the counted proof has to make is carried out inside the counted span, and no part of that deduction is delegated to this pointer.\par\medskip
\ifdefined\hypertarget\hypertarget{sourcepointer-2}{}\fi
\noindent\textit{Source reference.} Original lemma 02LO: sources/stacks-more-morphisms/pinned-original-source.tex, source lines 12136--12136. Navigation pointer only: the counted proof cites this statement by \textbackslash{}ref as a standing input; the deduction the counted proof has to make is carried out inside the counted span, and no part of that deduction is delegated to this pointer.\par\medskip
\begin{lemma}
\label{lemma-count-geometric-fibres-universally-open}
Let $f : X \to Y$ be a separated, locally quasi-finite, and universally open
morphism of schemes. Let $n_{X/Y}$ be as in
Lemma \ref{lemma-count-geometric-fibres}.
If $n_{X/Y}(y) \geq d$ for some $y \in Y$ and $d \geq 0$,
then $n_{X/Y} \geq d$ in an open neighbourhood of $y$.
\end{lemma}

\begin{proof}
The question is local on $Y$ hence we may assume $Y$ affine.
Let $K$ be an algebraic closure of the residue field $\kappa(y)$.
Our assumption is that $(X_y)_K$ has $\geq d$ connected components.
Then for a suitable quasi-compact open $X' \subset X$
the scheme $(X'_y)_K$ has $\geq d$ connected components; details omitted.
After replacing $X$ by $X'$ we may assume $X$ is quasi-compact.
Then $f$ is quasi-finite. Let $x_1, \ldots, x_n$ be the points of $X$
lying over $y$. Apply
Lemma \ref{lemma-etale-splits-off-quasi-finite-part-technical-variant}
to get an \'etale neighbourhood $(U, u) \to (Y, y)$ and a decomposition
$$
U \times_Y X =
W \amalg
\ \coprod\nolimits_{i = 1, \ldots, n}
\ \coprod\nolimits_{j = 1, \ldots, m_i}
V_{i, j}
$$
as in locus citatus. Observe that $n_{X/Y}(y) = \sum_i m_i$
in this situation; some details omitted.
Since $f$ is universally open, we see
that $V_{i, j} \to U$ is open for all $i, j$. Hence after shrinking
$U$ we may assume $V_{i, j} \to U$ is surjective
for all $i, j$. This proves that
$n_{U \times_Y X/U} \geq \sum_i m_i = n_{X/Y}(y) \geq d$.
Since the construction of $n_{X/Y}$ is compatible with
base change the proof is complete.
\end{proof}

\end{document}
